\documentclass[aspectratio=169,11pt]{beamer}
\usepackage[spanish,es-nodecimaldot,es-noshorthands]{babel}
\usepackage{booktabs,tikz}
\graphicspath{{resultados/figuras/}}
\input{resultados/tablas/reporte-cifras.tex}
\usetheme{default}
\setbeamercolor{structure}{fg=blue!55!black}
\setbeamertemplate{navigation symbols}{}
\setbeamertemplate{footline}[frame number]
\title{Asistencia gravitacional de Júpiter a Saturno}
\subtitle{Reto III.3 de Mecánica Clásica}
\author[Santiago Delgado y Angel Camilo Hincapie]{Santiago Delgado Cobos (2241400)\texorpdfstring{\\}{ y }Angel Camilo Hincapie Quintero (2241945)}
\institute{Universidad Industrial de Santander}
\date{hincapiec837@gmail.com}
\begin{document}
\frame{\titlepage}

\begin{frame}{Qué significan K, N, P y P0}
\small
\begin{tabular}{@{}p{.07\linewidth}p{.43\linewidth}p{.41\linewidth}@{}}
\toprule
 & Qué se integra & Función en el reto\\
\midrule
K & Kepler: un solo atractor en cada región. & Aproximación por cónicas empalmadas.\\[3pt]
N & Newton: Sol, Júpiter y Saturno a la vez, con posiciones de Horizons. & Referencia para medir el error de K.\\[3pt]
P & Fuerzas de N. Se traslada el origen dentro de las esferas. & Comprueba el efecto numérico de cambiar de marco.\\[3pt]
$P_0$ & Fuerzas de N fuera. Planeta y marea solar dentro, como en el cuaderno original. & Mide el efecto de una simplificación intermedia.\\
\bottomrule
\end{tabular}
\medskip

K y N bastan para la comparación principal. P comprueba su interpretación y P0 añade un diagnóstico. La $k$ minúscula del original es $GMm$. Al trabajar por unidad de masa se usa $\mu=GM$.
\end{frame}

\begin{frame}{De dónde salen los parámetros}
\small
\begin{tabular}{@{}p{.25\linewidth}p{.70\linewidth}@{}}
\toprule
Dato o parámetro & Origen y para qué sirve\\
\midrule
$\mu_i$, $\mathbf R_i$, $\dot{\mathbf R}_i$ & Los parámetros gravitacionales fijan las fuerzas. Las efemérides dan la posición y velocidad de los planetas para ubicarlos y cambiar de marco.\\[3pt]
$r_{\rm SOI}$ y $f$ & El radio de Laplace fija el cruce en K. $f=1$ es el nominal. Variar $f$ sirve para medir sensibilidad.\\[3pt]
$r_0,r_1,r_2$ & La altura de 500 km fija $r_0$. Los radios orbitales circulares $r_1,r_2$ fijan la elipse de Hohmann.\\[3pt]
$v_\infty,r_p,\delta,b$ & $v_\infty$ sale de la velocidad relativa. El periapsis $r_p$ se elige en el diseño y se calcula en K. El giro $\delta$ determina la salida. $b$ describe la geometría.\\[3pt]
Fase y umbral & Se ajusta la fase circular para llegar a Saturno. El umbral de 1000 km sirve solo para evaluar la diferencia entre los modelos.\\
\bottomrule
\end{tabular}
\end{frame}

\begin{frame}{Hohmann fija el estado de llegada a Júpiter}
\begin{columns}[T]
\column{.44\textwidth}
\centering\includegraphics[width=\linewidth,height=.57\textheight,keepaspectratio]{01_hohmann.pdf}
\column{.53\textwidth}
Perihelio en la Tierra, afelio en Júpiter. Sus radios fijan la elipse.
\[
a_t=\frac{r_1+r_2}{2},\quad
v_t^2=\mu_\odot\left(\frac2r-\frac1{a_t}\right).
\]
Vuelo de \DiasHohmann{} días. Júpiter parte $97{,}16^\circ$ por delante.
\[
v_{\infty,\oplus}=v_{t1}-v_1=8{,}793\;\text{km/s}.
\]
Desde la órbita baja de 500 km, $v_0$ es la rapidez circular y $\Delta v_0=\sqrt{v_{\infty,\oplus}^2+2v_0^2}-v_0=\ImpulsoLEO$ km/s.
\end{columns}
\vfill
{\scriptsize Ecuaciones originales de elipse, vis-viva y escape terrestre en \texttt{hohmann\_jupiter.tex}.}
\end{frame}

\begin{frame}{La hipérbola joviana determina el giro}
La sonda llega a $7{,}413$ km/s y Júpiter se mueve a $13{,}057$ km/s, ambos en sentido prógrado. La entrada relativa es retrógrada y $v_{\infty,J}=\ExcesoJoviano$ km/s.
\[
\varepsilon_J=\tfrac12v_{\infty,J}^2,\qquad
a_h=\frac{\mu_J}{v_{\infty,J}^2},\qquad |\mathbf h_J|=b\,v_{\infty,J}.
\]
\[
r_{p,J}=a_h(e_J-1),\qquad b=a_h\sqrt{e_J^2-1},\qquad
\delta=2\arcsin(1/e_J).
\]
Elegimos $r_{p,J}=10R_J$. Resultan $e_J=1{,}180$, $b=2{,}49$ millones de km y $\delta=\GiroJoviano^\circ$. $b$ mide la distancia perpendicular desde el centro a la asíntota entrante. Un mayor periapsis reduce el giro.
\vfill
{\scriptsize Relaciones de la hipérbola del original \texttt{hohmann\_jupiter.tex}.}
\end{frame}

\begin{frame}{La rotación cambia la energía respecto al Sol}
\begin{columns}[T]
\column{.53\textwidth}
Componentes radial y tangencial, en km/s.
\[
\begin{aligned}
\mathbf u_{\rm ent}&=(0,-5{,}643),\\
\mathbf u_{\rm sal}&=(5{,}076,2{,}466),\\
\mathbf V_J+\mathbf u_{\rm sal}&=(5{,}076,15{,}523).
\end{aligned}
\]
Las velocidades relativas tienen igual módulo. La energía solar aumenta en
\[
\Delta\varepsilon_\odot=\mathbf V_J\cdot
(\mathbf u_{\rm sal}-\mathbf u_{\rm ent})
=\EnergiaGanada\;\text{km}^2/\text{s}^2.
\]
\column{.43\textwidth}
\centering
\begin{tikzpicture}[scale=.67,>=stealth]
\draw[->,gray] (-.55,0)--(4.4,0) node[right,black] {$\hat{\mathbf r}$};
\draw[->,gray] (0,-2.5)--(0,3.1) node[above,black] {$\hat{\boldsymbol\theta}$};
\draw[->,very thick,blue!75!black] (0,0)--(0,-1.9) node[below left] {$\mathbf u_{\rm ent}$};
\draw[->,very thick,red!75!black] (0,0)--(1.71,.83) node[right] {$\mathbf u_{\rm sal}$};
\draw[->,red!75!black] (0,-.65) arc (-90:26:.65);
\node at (1,-.65) {$\delta$};
\end{tikzpicture}
\par\small Vectores relativos a Júpiter. El giro elegido da una componente radial saliente.
\end{columns}
\end{frame}

\begin{frame}{Del giro instantáneo al sobrevuelo de Saturno}
\begin{columns}[T]
\column{.54\textwidth}
La construcción asintótica aplica el giro en la posición orbital de Júpiter y en un solo instante. Desde ese estado se integran Sol y Saturno.
\medskip

Se ajusta la fase de Saturno a $24{,}92^\circ$ al paso por Júpiter. El mínimo es de \MinimoDiseno{} km tras \DiasDiseno{} días de vuelo desde Júpiter.
\column{.43\textwidth}
\centering\includegraphics[width=\linewidth,height=.55\textheight,keepaspectratio]{10_sobrevuelo_saturno.pdf}
\par\small Aproximación y salida en coordenadas centradas en Saturno.
\end{columns}
\end{frame}

\begin{frame}{Esferas: el encuentro se integra entre dos cruces}
\begin{columns}[T]
\column{.50\textwidth}
\centering
\begin{tikzpicture}[scale=.70,>=stealth]
\draw[dashed,blue!60!black] (4,0) circle (1.3);
\node[font=\scriptsize,blue!60!black] at (4,1.6) {$|\boldsymbol\rho_J|=r_{\rm SOI}$};
\fill (4,0) circle (.06) node[below right,font=\scriptsize]{Júpiter};
\draw[very thick,gray] (.2,-2.6) .. controls (1.8,-2.2) .. (3.08,-.92);
\draw[very thick,blue!70!black] (3.08,-.92) .. controls (3.4,-.2) and (3.9,.7) .. (4.65,1.126);
\draw[very thick,gray,->] (4.65,1.126)--(5.85,2.2);
\fill (3.08,-.92) circle (.055) node[below,font=\scriptsize]{entrada $t_{\rm in}$};
\fill (4.65,1.126) circle (.055) node[right,font=\scriptsize]{salida $t_{\rm out}$};
\node[font=\scriptsize] at (1.1,-1.55) {tramo solar};
\end{tikzpicture}
\par\small Esquema de cruces adaptado de \texttt{esferas\_influencia.tex}.
\column{.47\textwidth}
\[
r_{\rm SOI,i}(t)=f|\mathbf R_i(t)|
\left(\frac{\mu_i}{\mu_\odot}\right)^{2/5}.
\]
\small Entrada. Se restan $\mathbf R_i$ y $\dot{\mathbf R}_i$ al estado solar. Con el estado relativo se integra la cónica del planeta.
\medskip

Salida. Se suman $\mathbf R_i$ y $\dot{\mathbf R}_i$ en la fecha del cruce saliente. El estado físico es continuo.
\end{columns}
\end{frame}

\begin{frame}{Los lagrangianos del cuaderno y las tres variantes}
\small
Fuera de las esferas, el cuaderno añade a la cónica solar los potenciales directo e indirecto de cada planeta:
\[
\mathcal L_\odot=\tfrac12|\dot{\mathbf r}|^2+\frac{\mu_\odot}{r}
+\sum_{i=J,S}\mu_i\left[\frac1{|\mathbf r-\mathbf R_i|}
-\frac{\mathbf r\cdot\mathbf R_i}{R_i^3}\right].
\]
Dentro de la esfera, el corchete planetocéntrico genera la marea solar:
\[
\mathcal L_i^{(P0)}=\tfrac12|\mathbf u|^2+\frac{\mu_i}{\rho_i}
+\mu_\odot\left[\frac1{|\mathbf R_i+\boldsymbol\rho_i|}
+\frac{\mathbf R_i\cdot\boldsymbol\rho_i}{R_i^3}\right].
\]
\begin{itemize}
\item K quita los corchetes. Queda un atractor por región.
\item P0 retiene los corchetes del cuaderno y cambia de marco en la esfera.
\item N integra la ecuación de $\mathcal L_\odot$ en el marco solar. P conserva su dinámica al trasladarla: $\ddot{\boldsymbol\rho}_i=\ddot{\mathbf r}-\ddot{\mathbf R}_i$.
\end{itemize}
\end{frame}

\begin{frame}{Qué fija cada modelo y qué se compara}
\small
\begin{tabular}{@{}p{.18\linewidth}p{.37\linewidth}p{.36\linewidth}@{}}
\toprule
& Hohmann + dispersión & Esferas de influencia (K)\\
\midrule
Antes de Júpiter & Elipse tangente fija $v_{\infty,J}$. & El cruce entrante fija posición y velocidad finitas.\\[3pt]
En Júpiter & Se elige $r_{p,J}$ y se gira $\mathbf u$ en un instante. & $r_{p,J}$ y la duración resultan de integrar la cónica local.\\[3pt]
Hacia Saturno & El diseño propaga Sol y Saturno a partir del giro. & Sol entre esferas. Saturno actúa al entrar en su esfera.\\
\bottomrule
\end{tabular}
\medskip

Para medir el error de K, copiamos el estado de N en cada entrada joviana. Con $f=1$ la esfera tiene $49{,}20$ millones de km al entrar. K tarda unos 132 días en salir y queda fuera de la esfera de Saturno.
\end{frame}

\begin{frame}{El error depende del radio de empalme}
\centering
\begin{tabular}{rrrr}
\toprule
Factor $f$ & P (km) & P0 (km) & K (km)\\
\midrule
\input{resultados/tablas/reporte-frontera.tex}
\end{tabular}
\par\medskip
\raggedright
Diferencias respecto a N. P0 conserva la fuerza del planeta y la marea solar dentro de cada esfera. Con un umbral de 1000 km, P cumple en todos los casos. Las cónicas con radio nominal pasan fuera de la esfera de Saturno.
\end{frame}

\begin{frame}{Qué respaldan las verificaciones}
\begin{itemize}
\item Una integración independiente reproduce la dispersión hiperbólica analítica. En K, la energía jovicéntrica varía numéricamente del orden de $10^{-10}$ en términos relativos.
\item Al reducir el paso y las tolerancias, el residuo en posición P--N al final baja de 7045 a 6,30 km.
\item NASA publica 161\,081 km para Voyager 2 en Saturno. La interpolación diaria da 151\,408 km y N predice 4,34 millones de km.
\end{itemize}
\medskip
Cerca del periapsis de Júpiter domina la fuerza central. El error de omitir perturbaciones cerca de la frontera se amplifica en el viaje a Saturno.
\end{frame}

\begin{frame}{Diseño con efemérides y alcance}
Ajustamos K y N por separado para apuntar a un lugar desplazado un millón de km del centro de Saturno. Los mínimos resultantes son de unos 982\,452 y 983\,201 km, respectivamente. Sus épocas difieren en casi 16 días.
\medskip

La prueba con el mismo estado de entrada mide el error de la aproximación. El ajuste independiente muestra que ambos modelos permiten diseñar un sobrevuelo bajo sus propias hipótesis.
\medskip

\footnotesize Fuentes. Efemérides JPL Horizons y distancias históricas NASA NSSDCA. El repositorio contiene código, datos y pruebas. La bitácora documenta las intervenciones de IA.
\end{frame}
\end{document}
